% Part: first-order-logic
% Chapter: syntax-and-semantics
% Section: semantic-notions

\documentclass[../../../include/open-logic-section]{subfiles}

\begin{document}

\olfileid{qml}{syn}{sem}

\olsection{Semantic Notions}

\begin{explain}
The notions of logical truth, validity, satisfiability, are
defined as in modal logic. This holds for free quantified 
modal logic as well as classical quantified modal logic.

A logical truth $!A$ is a sentence that is true at every world of every model.

An argument with set of sentences $\Gamma$ as premises and sentence
$!A$ as conclusion is valid (written $\Gamma \Entails !A$) iff 
there is no model and world at which all the sentences in $\Gamma$ 
are true but $!A$ is false. 

\Article{sentence} !!{sentence} is satisfiable iff there is a model and a world 
at which it is true.

\Article{formula} !!{formula} is satisfiable iff it satisfied at some
world relative to some assignment in some model. 

And so on.

\end{explain}

\end{document}
